\documentclass[12pt,a4paper]{article}

% ====== ЯЗЫК И ШРИФТЫ (XeLaTeX) ======
\usepackage{fontspec}
\usepackage{polyglossia}
\setdefaultlanguage{russian}
\setotherlanguage{english}

\setmainfont{CMU Serif}
\setsansfont{CMU Sans Serif}
\setmonofont{CMU Typewriter Text}
\newfontfamily\cyrillicfont{CMU Serif}
\newfontfamily\cyrillicfontsf{CMU Sans Serif}
\newfontfamily\cyrillicfonttt{CMU Typewriter Text}

\usepackage{microtype}


% ====== МАТЕМАТИКА ======
\usepackage{amsmath, amssymb, amsthm, mathtools}
\usepackage{bm} % Жирные греческие буквы


% Русские стандарты математической типографики
\renewcommand{\ge}{\geqslant}
\renewcommand{\le}{\leqslant}
\newcommand{\E}{\mathbb{E}}
\renewcommand{\P}{\mathbb{P}}
\newcommand{\F}{\mathcal{F}}

% ====== ГРАФИКА И TIKZ ======
\usepackage{tikz}
\usetikzlibrary{arrows.meta, positioning, automata, backgrounds, decorations.pathmorphing}

% ====== ОФОРМЛЕНИЕ СТРАНИЦЫ ======
\usepackage[left=2cm, right=2cm, top=2.5cm, bottom=2.5cm]{geometry}
\usepackage{parskip} % Красивые отступы между абзацами
\usepackage{xcolor}


% ====== ГИПЕРССЫЛКИ ======
\usepackage[colorlinks=true, linkcolor=blue!70!black, urlcolor=blue!70!black, citecolor=green!60!black]{hyperref}

% ====== TCOLORBOX (ОКРУЖЕНИЯ ДЛЯ ТЕОРЕМ И ЗАДАЧ) ======
\usepackage[most]{tcolorbox}
\tcbuselibrary{skins, breakable, theorems}

% Цветовая палитра "МЦНМО/ВШЭ"
\definecolor{defblue}{RGB}{20, 75, 140}
\definecolor{thmred}{RGB}{160, 30, 45}
\definecolor{probgreen}{RGB}{25, 110, 55}
\definecolor{warnpurple}{RGB}{130, 30, 130}
\definecolor{insightteal}{RGB}{0, 100, 100}
\definecolor{solgray}{RGB}{90, 90, 90}

% Определение
\newtcbtheorem[auto counter, number within=section]{definition}{Определение}{
    enhanced, breakable,
    colback=defblue!3, colframe=defblue,
    fonttitle=\bfseries\sffamily, title style={color=defblue},
    attach boxed title to top left={yshift=-2.5mm, xshift=5mm},
    boxed title style={sharp corners},
    boxrule=1pt, drop fuzzy shadow
}{def}

% Теорема
\newtcbtheorem[auto counter, number within=section]{theorem}{Теорема}{
    enhanced, breakable,
    colback=thmred!3, colframe=thmred,
    fonttitle=\bfseries\sffamily, title style={color=thmred},
    attach boxed title to top left={yshift=-2.5mm, xshift=5mm},
    boxed title style={sharp corners},
    boxrule=1pt, drop fuzzy shadow
}{thm}

% Задача
\newtcbtheorem[auto counter, number within=section]{problem}{Задача}{
    enhanced, breakable,
    colback=probgreen!3, colframe=probgreen,
    fonttitle=\bfseries\sffamily, title style={color=probgreen},
    attach boxed title to top left={yshift=-2.5mm, xshift=5mm},
    boxed title style={sharp corners},
    boxrule=1pt, drop fuzzy shadow
}{prob}

% Важное замечание / Парадокс
\newtcolorbox{paradox}[1][]{
    enhanced, breakable,
    colback=warnpurple!4, colframe=warnpurple!90!black,
    title={\textbf{Внимание: #1}},
    fonttitle=\bfseries\sffamily,
    coltitle=black,
    attach title to upper,
    after title={\par\smallskip},
    boxrule=1.5pt, arc=4pt
}

% Олимпиадный инсайт
\newtcolorbox{insight}[1][]{
    enhanced, breakable,
    colback=insightteal!3, colframe=insightteal, leftrule=4pt,
    boxrule=0pt, arc=0pt,
    title={\textbf{Олимпиадный инсайт: #1}},
    coltitle=insightteal, fonttitle=\bfseries\sffamily,
    attach title to upper,
    after title={\par\smallskip}
}

% Решение
\newenvironment{solution}{
    \begin{tcolorbox}[
        enhanced, breakable, blanker,
        borderline west={2.5pt}{0pt}{solgray},
        left=12pt, top=5pt, bottom=5pt,
        before skip=10pt, after skip=15pt
    ]
    \textbf{\textsf{РЕШЕНИЕ.}}\;
}{
    \hfill $\blacksquare$
    \end{tcolorbox}
}



\begin{document}

\begin{center}
    \vspace*{1.0cm}
    {\Huge \bfseries Корни Времени\par}
    \vspace{0.4cm}
    {\Large \itshape Искусство рекуррентных и характеристических уравнений\par}
    \vspace{0.6cm}
    \rule{0.6\textwidth}{1pt}\par
    \vspace{0.6cm}
    {\large \textbf{Калашников Александр}}\par
    \vspace{0.2cm}
    {\small Wildberries \& Russ, ФКН ВШЭ}
    \vspace{1.0cm}
\end{center}

\tableofcontents
\newpage

\section{Введение: Предел возможностей \texttt{for}}

В олимпиадном программировании (ДП) и теории графов мы постоянно описываем состояния системы шаг за шагом: количество путей $x_n$ зависит от $x_{n-1}$ и $x_{n-2}$. Компьютер решает это циклом \texttt{for} за $\mathcal{O}(n)$. 

Но в олимпиадной математике нас часто просят найти \textbf{точную аналитическую формулу} $n$-го члена, или найти предел при $n \to \infty$. Более того, если $n = 10^{18}$, любой цикл получит \texttt{Time Limit Exceeded}. На сцену выходит теория, которая превращает бесконечную цепочку локальных зависимостей в замкнутую глобальную формулу, вычисляемую за $\mathcal{O}(1)$.

\section{Фундамент: Откуда берётся $\lambda^n$?}

\begin{definition}{Линейная рекуррента $k$-го порядка}{lin_rec}
Последовательность $\{x_n\}$ называется линейной однородной рекуррентой с постоянными коэффициентами, если для любого $n \ge k$:
\[ x_n = a_1 x_{n-1} + a_2 x_{n-2} + \dots + a_k x_{n-k} \]
\end{definition}

Главный метод решения — так называемый \textit{Анзац Эйлера}: мы ищем решение в виде $x_n = \lambda^n$ ($\lambda \neq 0$).

\begin{insight}[Почему мы ищем решение именно в виде экспоненты?]
В математическом анализе есть оператор взятия производной $D(f) = f'$. Его \textbf{собственными функциями} (теми, которые при применении оператора не меняют форму, а лишь умножаются на константу) являются экспоненты: $D(e^{\lambda t}) = \lambda \cdot e^{\lambda t}$. Именно поэтому дифференциальные уравнения решаются через подстановку $e^{\lambda t}$.

В дискретной математике аналогом производной является \textbf{оператор сдвига} $S(x_n) = x_{n+1}$. Какая дискретная функция при сдвиге индекса на $+1$ просто умножается на константу? Только геометрическая прогрессия! 
$S(\lambda^n) = \lambda^{n+1} = \lambda \cdot \lambda^n$. 
Подставляя $x_n = \lambda^n$, мы переводим сложный сдвиг во времени в обычное алгебраическое уравнение.
\end{insight}

Подставим $x_n = \lambda^n$ в рекурренту 2-го порядка $x_{n+2} = a x_{n+1} + b x_n$:
\[ \lambda^{n+2} = a \lambda^{n+1} + b \lambda^n \]
Разделим обе части на $\lambda^n$ и перенесем всё влево. Получаем квадратное уравнение:

\begin{theorem}{Характеристическое уравнение}{char_eq}
Рекурренте $x_{n+2} - a x_{n+1} - b x_n = 0$ ставится в соответствие \textbf{характеристическое уравнение}:
\[ \lambda^2 - a \lambda - b = 0 \]
Структура общего решения зависит от корней этого уравнения:
\begin{itemize}
    \item \textbf{Корни различны ($\lambda_1 \neq \lambda_2 \in \mathbb{R}$):} Базис двумерен. Общее решение — линейная комбинация:
    \[ x_n = C_1 \cdot \lambda_1^n + C_2 \cdot \lambda_2^n \]
    \item \textbf{Корень кратный ($\lambda_1 = \lambda_2 = \lambda$):} Возникает резонанс. Второе решение домножают на $n$:
    \[ x_n = C_1 \cdot \lambda^n + C_2 \cdot \mathbf{n} \cdot \lambda^n \]
    \item \textbf{Комплексные корни ($D < 0$):} Корни имеют вид $r(\cos \varphi \pm i \sin \varphi)$. Решение осциллирует:
    \[ x_n = r^n \left( C_1 \cos(\varphi n) + C_2 \sin(\varphi n) \right) \]
\end{itemize}
Константы $C_1$ и $C_2$ однозначно определяются из начальных условий (СЛАУ размера $2 \times 2$).
\end{theorem}

\newpage
\section{Аномалии: Комплексные корни и Тайна Резонанса}

Школьники часто теряются, когда алгоритм ломается (дискриминант $<0$ или $D=0$). Давайте посмотрим на шестерёнки под капотом этих явлений.

\subsection{Тайна кратных корней: откуда берется «$n$»?}

\begin{paradox}[Эффект производной]
В школе говорят: "Если корни совпали, $\lambda_1 = \lambda_2 = \lambda$, то вторая независимая функция — это $n \lambda^n$". Но откуда взялось это $n$? Почему не $n^2$ и не $\log n$?
\end{paradox}

Рассмотрим левую часть рекурренты как оператор: $P(\lambda) = \lambda^n - a\lambda^{n-1} - b\lambda^{n-2} = 0$. 
Из алгебры известно, что если многочлен имеет корень $\lambda_0$ кратности 2, то в этой точке равен нулю не только сам многочлен, но и \textbf{его производная по параметру $\lambda$}: $P'(\lambda_0) = 0$.

Давайте продифференцируем наше уравнение $P(\lambda)=0$ по переменной $\lambda$:
\[ \frac{\partial}{\partial \lambda} \left( \lambda^n - a\lambda^{n-1} - b\lambda^{n-2} \right) = n\lambda^{n-1} - a(n-1)\lambda^{n-2} - b(n-2)\lambda^{n-3} = 0 \]
Домножим обе части на $\lambda$:
\[ \bm{n\lambda^n} - a\bm{(n-1)\lambda^{n-1}} - b\bm{(n-2)\lambda^{n-2}} = 0 \]
Присмотритесь! Это выражение в точности означает, что новая последовательность $y_n = n \lambda^n$ также удовлетворяет исходной рекурренте $y_n - a y_{n-1} - b y_{n-2} = 0$. 
Множитель $n$ возникает прямо из определения производной $\frac{d}{d\lambda}(\lambda^n) = n \lambda^{n-1}$! Мы восстановили потерянную размерность векторного пространства.

\subsection{Осцилляции и Тригонометрия ($D < 0$)}

Что делать, если вещественных корней нет? По формуле Эйлера $e^{i\varphi} = \cos \varphi + i \sin \varphi$. Если характеристическое уравнение имеет комплексные корни $\lambda = r e^{\pm i\varphi}$, это означает, что система ведёт себя как маятник — она \textbf{периодически колеблется}!

\begin{problem}{Осциллирующий сервер (Классика ВсОШ)}{complex}
Нагрузка на сервер $a_n$ меняется по закону $a_n = a_{n-1} - a_{n-2}$. Известно, что в 1-ю секунду нагрузка была 10, а во 2-ю — 20. Какой она будет на 2026-ю секунду?
\end{problem}
\begin{solution}
Характеристическое уравнение: $\lambda^2 - \lambda + 1 = 0$. 
Дискриминант $D = 1 - 4 = -3 < 0$. Корни: $\lambda_{1,2} = \frac{1 \pm i\sqrt{3}}{2}$. (Целочисленная рекуррента породила мнимые числа!).

Переведем корень в тригонометрическую форму. Модуль $r = \sqrt{(1/2)^2 + (\sqrt{3}/2)^2} = 1$. Угол $\varphi$: $\cos \varphi = 1/2 \implies \varphi = \pi/3$ (или $60^\circ$).
Так как $r=1$, решение состоит только из чистых синусов и косинусов:
\[ a_n = C_1 \cos\left(\frac{\pi n}{3}\right) + C_2 \sin\left(\frac{\pi n}{3}\right) \]

\textbf{Олимпиадный хак:} Нам даже не нужно искать $C_1$ и $C_2$! Мы видим, что аргумент функций равен $\frac{\pi n}{3}$. Поскольку период синуса и косинуса равен $2\pi$, функция полностью повторит свои значения, когда $\frac{\pi n}{3} = 2\pi \implies \mathbf{n = 6}$.
Последовательность зациклена с периодом 6. 
Остаток от деления 2026 на 6 равен 4. Мы знаем $a_1=10, a_2=20$. Найдём $a_4$ вручную:
$a_3 = a_2 - a_1 = 20 - 10 = 10$. \\
$a_4 = a_3 - a_2 = 10 - 20 = -10$. \\
Ответ: $\mathbf{-10}$.
\end{solution}

\newpage
\section{Комбинаторные графы и Системы}

Часто система описывается несколькими переменными. В этом случае нужно составить систему уравнений и свести её к одной переменной методом подстановки.

\begin{problem}{Маршрутизация дата-пакета}{packet}
Сетевой пакет маршрутизируется в кольцевой топологии из трех серверов: A, B и C. Изначально пакет находится на сервере C. Затем он начинает пересылаться между серверами, каждый раз переходя на один из двух других (соседних). \\
Сколько существует различных маршрутов для пакета, чтобы, совершив ровно 239 пересылок, он в итоге снова оказался на сервере C?
\end{problem}

\begin{solution}
Обозначим количество путей длины $n$, оканчивающихся в C, через $c_n$.
В силу симметрии кольца, количество путей в A равно количеству путей в B. Обозначим это за $a_n$.

Выпишем систему переходов (откуда мы могли прийти на текущем шаге):
\begin{align*}
    c_n &= a_{n-1} + b_{n-1} = 2 a_{n-1} \\
    a_n &= b_{n-1} + c_{n-1} = a_{n-1} + c_{n-1}
\end{align*}

Из первого уравнения $a_{n-1} = \frac{1}{2}c_n$. Сдвинем индекс на $+1$: $a_n = \frac{1}{2}c_{n+1}$. 
Подставим это во второе уравнение:
\[ \frac{1}{2}c_{n+1} = \frac{1}{2}c_n + c_{n-1} \implies \mathbf{c_{n+1} - c_n - 2c_{n-1} = 0} \]

Получили рекурренту! Характеристическое уравнение: $\lambda^2 - \lambda - 2 = 0$.
По теореме Виета корни: $\lambda_1 = 2$, $\lambda_2 = -1$.
Общий вид решения: $c_n = C_1 \cdot 2^n + C_2 \cdot (-1)^n$.

Зададим начальные условия. За $0$ шагов пакет УЖЕ в C: $c_0 = 1$. За $1$ шаг он обязан уйти: $c_1 = 0$.
\begin{cases}
C_1 + C_2 = 1 \\
2C_1 - C_2 = 0 \implies C_2 = 2C_1
\end{cases}
Отсюда $3C_1 = 1 \implies C_1 = 1/3, \quad C_2 = 2/3$.

Формула для любого $n$: $c_n = \frac{2^n + 2 \cdot (-1)^n}{3}$. \\
Для нечётного $n=239$ получаем $(-1)^{239} = -1$. Ответ: $\mathbf{\frac{2^{239} - 2}{3}}$.
\end{solution}

\section{Столкновение с вероятностью (Марковские цепи)}

В задачах на цепи Маркова рекурренты возникают из \textit{формулы полной вероятности}.

\begin{problem}{Стохастический парсер}{parser}
Стохастический генератор на каждом такте равновероятно выдает число 1 или 2. Скрипт запускает генератор много раз подряд, прибавляя число к общей сумме (изначально 0). \\
Найдите вероятность того, что в какой-то момент общая сумма станет в точности равна 2024.
\end{problem}

\begin{solution}
Пусть $p_n$ — вероятность того, что общая сумма когда-нибудь \textit{наступит} на значение $n$.
Чтобы получить сумму $n$, мы должны были непосредственно перед этим либо:
\begin{itemize}
    \item Оказаться в сумме $(n-1)$ и выбросить «1» (вероятность $\frac{1}{2}$).
    \item Оказаться в сумме $(n-2)$ и выбросить «2» (вероятность $\frac{1}{2}$).
\end{itemize}
По формуле полной вероятности:
\[ p_n = \frac{1}{2}p_{n-1} + \frac{1}{2}p_{n-2} \implies \mathbf{2p_n - p_{n-1} - p_{n-2} = 0} \]

Характеристическое уравнение: $2\lambda^2 - \lambda - 1 = 0$. Корни: $\lambda_1 = 1, \; \lambda_2 = -\frac{1}{2}$.
Общий вид: $p_n = C_1 \cdot 1^n + C_2 \cdot \left(-\frac{1}{2}\right)^n = C_1 + C_2 \left(-\frac{1}{2}\right)^n$.

База: $p_0 = 1$ (мы стартуем из нуля) и $p_1 = \frac{1}{2}$ (нужно кинуть 1 первым броском).
Решая СЛАУ для $C_1, C_2$, получаем $C_1 = \frac{2}{3}, \quad C_2 = \frac{1}{3}$.

Формула вероятности: $p_n = \frac{2}{3} + \frac{1}{3} \left(-\frac{1}{2}\right)^n$. \\
Для чётного $n=2024$ ответ: $\mathbf{\frac{2}{3} + \frac{1}{3 \cdot 2^{2024}}}$.
\end{solution}

\begin{insight}[Смысл Предела и Эргодичность]
При $n \to \infty$ слагаемое $(-1/2)^n$ стремится к нулю, и $p_n \to \frac{2}{3}$. Это потрясающе логично физически: матожидание одного броска генератора равно $\frac{1+2}{2} = 1.5$. В среднем за один такт мы проходим расстояние 1.5. Значит, на длинной прямой мы "наступаем" на $\frac{1}{1.5} = \frac{2}{3}$ всех целых чисел! Наша рекуррента вывела \textit{Теорему Восстановления} для случайных процессов чистой алгеброй.
\end{insight}

\newpage
\section{Изнанка Мартингалов: Тайна Игры с Разорением}

В методичке по Мартингалам мы решали задачу о нечестном казино. Игрок выигрывает $1$ рубль с вероятностью $p$ и проигрывает с вероятностью $q$. Чтобы найти вероятность достижения цели $N$, мы использовали инвариант Муавра $\left(\frac{q}{p}\right)^x$. 

Школьники часто спрашивают: \textit{«Как догадаться до этой функции без магии?»} Отвечают характеристические уравнения!

\begin{problem}{Вероятность разорения (Gambler's Ruin)}{ruin}
У игрока $x$ рублей. Он хочет дойти до $N$ рублей. На каждом шаге он выигрывает 1 рубль с вероятностью $p$ и проигрывает с вероятностью $q$ ($p \neq q, p+q=1$). При $0$ он разоряется. Найдите вероятность успеха $P_x$.
\end{problem}

\begin{solution}
Посмотрим на \textbf{один шаг вперёд} из состояния $x$. По формуле полной вероятности:
\[ P_x = p \cdot P_{x+1} + q \cdot P_{x-1} \]
Перепишем это в виде классической рекурренты (выразив старший индекс):
\[ p \cdot P_{x+1} - P_x + q \cdot P_{x-1} = 0 \]

Составляем характеристическое уравнение (замена $P_x \to \lambda^x$):
\[ p\lambda^2 - \lambda + q = 0 \]
Так как $p+q=1$, изящно заменим коэффициент $-\lambda$ на $-(p+q)\lambda$:
\[ p\lambda^2 - p\lambda - q\lambda + q = 0 \implies p\lambda(\lambda - 1) - q(\lambda - 1) = 0 \implies (\lambda - 1)(p\lambda - q) = 0 \]
Корни уравнения: $\lambda_1 = 1$ и $\bm{\lambda_2 = \frac{q}{p}}$. 

Бинго! Вот откуда берётся экспоненциальный мартингал! 
Общее решение:
\[ P_x = C_1 \cdot 1^x + C_2 \cdot \left(\frac{q}{p}\right)^x = C_1 + C_2 \left(\frac{q}{p}\right)^x \]

Используя граничные условия (Базу) $P_0 = 0$ и $P_N = 1$, мы элементарно находим $C_1, C_2$ и получаем ту самую формулу $P_x = \frac{(q/p)^x - 1}{(q/p)^N - 1}$.
\end{solution}

\begin{paradox}[Секретная связь: Мартингалы и Рекурренты]
\textbf{Мартингалы} смотрят на процесс "сверху вниз": мы требуем сохранения системы в матожидании $\E[f(X_{n+1}) \mid X_n] = f(X_n)$, и пытаемся угадать функцию $f(x)$. \\
\textbf{Рекурренты} смотрят "снизу вверх": мы выписываем формулу перехода и честно находим базис. Оказывается, стохастические инварианты — это просто переодетые \textbf{корни характеристического уравнения}! Магия полностью разоблачена.
\end{paradox}

\newpage
\section{Неоднородные уравнения: Штраф за время}

Уравнение вида $x_n - a x_{n-1} - b x_{n-2} = f(n)$ называется \textbf{неоднородным}. 
С точки зрения физики, однородная часть — это свободные колебания системы, а $f(n)$ — вынуждающая внешняя сила. Общее решение состоит из двух частей:
\[ x_n = x_n^{(0)} + x_n^{(*)} \]
где $x_n^{(0)}$ --- общее решение однородной задачи ($= 0$), а $x_n^{(*)}$ --- \textbf{любое частное решение} неоднородной задачи, которое мы угадываем (Метод неопределенных коэффициентов).

\begin{tcolorbox}[colback=white, colframe=defblue!80!black, title=\textbf{Таблица Угадываний (Анзац для $x_n^{(*)}$)}, fonttitle=\bfseries]
\renewcommand{\arraystretch}{1.5}
\begin{tabular}{|p{4.5cm}|p{5.5cm}|p{5cm}|}
\hline
\textbf{Правая часть $f(n)$} & \textbf{Ищем $x_n^{(*)}$ в виде:} & \textbf{Условие резонанса ($n^s$)} \\
\hline
Константа $C$ & Константа $A$ & Умножить на $n^s$, если $\lambda=1$ — корень кратности $s$ \\
\hline
Многочлен $P_k(n)$ & Многочлен $A n^k + B n^{k-1} \dots$ & Умножить на $n^s$, если $\lambda=1$ — корень кратности $s$ \\
\hline
Экспонента $C \cdot d^n$ & Экспонента $A \cdot d^n$ & Умножить на $n^s$, если $\lambda=d$ — корень кратности $s$ \\
\hline
\end{tabular}
\end{tcolorbox}
\textbf{⚠️ Суть Резонанса:} Если внешняя сила бьёт в собственную частоту системы (уже содержится в однородном решении), возникает бесконечный линейный или квадратичный рост!

\begin{problem}{Происхождение квадратичного мартингала}{}
Игрок с капиталом $x$ играет в честную орлянку ($p=q=0.5$). Игра заканчивается при достижении $N$ или $0$. Сколько в среднем бросков она продлится?
\end{problem}
\begin{solution}
Пусть $E_x$ — матожидание оставшихся шагов. Тратим 1 шаг (штраф времени $+1$):
\[ E_x = \frac{1}{2} E_{x+1} + \frac{1}{2} E_{x-1} + \mathbf{1} \implies \mathbf{E_{x+1} - 2E_x + E_{x-1} = -2} \]

\textbf{1. Однородная часть:} $\lambda^2 - 2\lambda + 1 = 0 \implies \lambda = 1$ (кратность $s=2$). 
Однородное решение: $E_x^{(0)} = C_1 + C_2 x$.

\textbf{2. Частное решение:} Справа константа $f(n) = -2$. По таблице мы ищем ответ в виде константы $A$. НО! Корень $\lambda=1$ имеет кратность 2. Возникает \textbf{двойной резонанс}. Мы обязаны домножить наш Анзац на $x^2$. 
Ищем решение в виде параболы $E_x^{(*)} = A \cdot x^2$. Подставляем:
\[ A(x+1)^2 - 2Ax^2 + A(x-1)^2 = -2 \implies 2A = -2 \implies A = -1 \]
Частное решение: $E_x^{(*)} = \mathbf{-x^2}$.

\textbf{3. Склейка и база:} $E_x = C_1 + C_2 x - x^2$. Граничные условия $E_0 = 0, E_N = 0$.
Из них сразу следует $C_1 = 0$ и $C_2 = N$. Итоговая формула: $E_x = \mathbf{x(N - x)}$. 
Мы только что математически строго доказали инвариант времени $\mathbf{X_n^2 - n}$! 
\end{solution}

\begin{problem}{Ожидаемое время нечестной игры (Дрейф)}{}
А что если игра нечестная? $p \neq q$. Найти общее решение для ожидаемого времени.
\end{problem}
\begin{solution}
Уравнение: $p E_{x+1} - E_x + q E_{x-1} = -1$. Корни: $\lambda_1 = 1$ и $\lambda_2 = q/p$.
Так как $p \neq q$, корень $\lambda=1$ имеет кратность 1. Возникает \textbf{одинарный} резонанс! Ищем решение в виде прямой: $E_x^{(*)} = A \cdot x$. Подставляем:
\[ p A(x+1) - Ax + q A(x-1) = -1 \implies Ax(p-1+q) + A(p-q) = -1 \]
Так как $p+q=1$, $A(p - q) = -1 \implies A = \frac{1}{q-p}$.
Общее решение: $E_x = C_1 + C_2 \left(\frac{q}{p}\right)^x + \mathbf{\frac{x}{q-p}}$. 

\textit{Фантастический вывод: в нечестной игре время растет линейно $\mathcal{O}(x)$, а в честной — квадратично $\mathcal{O}(x^2)$. Знаменитое Тождество Вальда доказано нами чистой алгеброй!}
\end{solution}

\newpage
\section{Бонус: Матричное возведение и Теорема Гамильтона-Кэли}

На олимпиадах по программированию (и в Machine Learning, например, в RNN-сетях) часто просят найти ответ для гигантского $n=10^{18}$ по модулю $10^9+7$. Формула Бине с $\sqrt{5}$ даст огромную погрешность `float` на 20-м знаке.

Здесь применяется \textbf{Линейная Алгебра}. Любую рекурренту записывают в векторно-матричном виде $\vec{v}_n = M \vec{v}_{n-1}$. Для Фибоначчи:
\[
\begin{pmatrix} F_{n+1} \\ F_n \end{pmatrix} = 
\underbrace{\begin{pmatrix} 1 & 1 \\ 1 & 0 \end{pmatrix}}_{M}
\begin{pmatrix} F_n \\ F_{n-1} \end{pmatrix}
\]
Мы можем возвести матрицу $M$ в степень $n$ \textbf{бинарным возведением за $\mathcal{O}(\log n)$}, сохраняя 100\% точность целых чисел!

Но самое красивое впереди. Как найти \textbf{Собственные значения (eigenvalues)} матрицы переходов $M$? Из уравнения $\det(M - \lambda I) = 0$.
\[
\det \begin{pmatrix} 1-\lambda & 1 \\ 1 & -\lambda \end{pmatrix} = (1-\lambda)(-\lambda) - 1 = \lambda^2 - \lambda - 1
\]
ШОК! Характеристический многочлен матрицы \textbf{в точности совпадает} с характеристическим уравнением нашей рекурренты! Графы (матрица смежности), марковские цепи (матрица переходов) и рекурренты — это три разных языка для описания одной сущности: Линейного Оператора. И его поведение полностью диктуется его спектром (корнями).

\vspace{0.6cm}
\begin{tcolorbox}[colback=yellow!5, colframe=orange!80!black, title=\textbf{Шпаргалка Олимпиадника: Алгоритм решения}, fonttitle=\bfseries]
\begin{enumerate}
    \item \textbf{Каноничный вид:} Соберите все неизвестные с одной стороны $x_n - a_1 x_{n-1} \dots = f(n)$.
    \item \textbf{Корни однородной части ($f(n)=0$):} Замените $x_{n-i}$ на $\lambda^{k-i}$. Найдите корни $\lambda$.
    \item \textbf{Сформируйте базис $x_n^{(0)}$:}
    \begin{itemize}
        \item Разные вещественные корни: $C_1 \lambda_1^n + C_2 \lambda_2^n$.
        \item Кратные корни (вступает производная!): $(C_1 + C_2 \cdot n) \lambda^n$.
        \item Комплексные корни: Синусы/косинусы $r^n (A \cos \varphi n + B \sin \varphi n)$.
    \end{itemize}
    \item \textbf{Правая часть $x_n^{(*)}$ (Неоднородность):} Угадайте частное решение по таблице. \textbf{Помните про резонанс}: при совпадении с корнями домножайте на $n^k$.
    \item \textbf{База индукции:} В финальную сумму $x_n = x_n^{(0)} + x_n^{(*)}$ подставьте $n=0, 1 \dots$, решите СЛАУ.
\end{enumerate}
\end{tcolorbox}

\end{document}